\documentclass[12pt]{article}
\usepackage[margin=1in]{geometry}
\usepackage{amsmath,amssymb,amsfonts,amsthm,bm,array,booktabs,enumitem}
\usepackage[colorlinks=true,linkcolor=blue,citecolor=blue,urlcolor=blue]{hyperref}

\numberwithin{equation}{section}

\newtheorem{theorem}{Theorem}
\newtheorem{proposition}{Proposition}
\newtheorem{lemma}{Lemma}
\newtheorem{assumption}{Assumption}
\theoremstyle{definition}
\newtheorem{definition}{Definition}
\newtheorem{remark}{Remark}

\newcommand{\E}{\mathbb E}
\newcommand{\R}{\mathbb R}
\newcommand{\Pp}{\mathbb P}
\newcommand{\ind}{\mathbf 1}
\newcommand{\op}{o_p}
\newcommand{\Op}{O_p}
\newcommand{\N}{\mathcal N}
\newcommand{\Hh}{\mathcal H}
\newcommand{\LtwoT}{L^2(P_T)}
\newcommand{\LtwoP}{L^2(P)}
\newcommand{\innerT}[2]{\langle #1,#2\rangle_T}
\newcommand{\innerP}[2]{\langle #1,#2\rangle_P}

\begin{document}

\begin{center}
{\Large\bf Minimal Main Body for Profile Sieve EL}\\[6pt]
{\large Minimum Geometry Needed to Understand the Proof}\\[6pt]
{\normalsize Standalone main-body version}\\[12pt]
\end{center}

\tableofcontents
\newpage

\section{Model and Geometry}

The finite sample is viewed as the empirical Hilbert space
\[
    \LtwoT=(\R^T,\innerT{a}{b}=T^{-1}a'b),
    \qquad \|a\|_T^2=\innerT{a}{a}.
\]
The population limit is
\[
    \LtwoP,\qquad \innerP{f}{g}=\E\{fg\}.
\]
The sample model is
\[
    \bm y=\beta_0\bm x+\bm\mu_0+\bm u,
\]
where
\[
    \mu_{0,t}=m_{j_0(t)}(W_{t-1}),
    \qquad
    j_0(t)=j\iff k_{j,0}<t\le k_{j+1,0}.
\]
Thus \(\bm\mu_0\) is the stacked nuisance mean. If there are no nuisance
breaks, \(\mu_{0,t}=m(W_{t-1})\). With nuisance breaks, the function changes by
regime.

A candidate partition is
\[
    \Lambda=(k_1,\ldots,k_q),
    \qquad
    0=k_0<k_1<\cdots<k_q<k_{q+1}=T,
\]
with segments
\[
    I_j(\Lambda)=\{t:k_j<t\le k_{j+1}\}.
\]
Let \(P\) be the \(T\times K\) sieve matrix with row
\[
    P_K(W_{t-1})'=(p_1(W_{t-1}),\ldots,p_K(W_{t-1})).
\]
Let \(D_j(\Lambda)\) be the diagonal selector for segment \(I_j(\Lambda)\), and
define the block-sieve design
\[
    \bm Q_{K,\Lambda}
    =
    [D_0(\Lambda)P,\ldots,D_q(\Lambda)P].
\]
The candidate nuisance subspace is
\[
    \N_{K,\Lambda,T}
    =
    \operatorname{col}(\bm Q_{K,\Lambda})
    \subset \LtwoT.
\]
Let \(\Pi_{K,\Lambda,T}\) be the empirical orthogonal projection onto this
subspace and
\[
    \bm M_{K,\Lambda}=I-\Pi_{K,\Lambda,T}
\]
the residual-maker. The proof is governed by the subspace-valued map
\[
    \Lambda\longmapsto \N_{K,\Lambda,T}.
\]
For a trial \(\beta\), define
\[
    \bm y_\beta=\bm y-\beta\bm x.
\]
The profiled residual \(\bm M_{K,\Lambda}\bm y_\beta\) and the residualized
score direction \(\bm M_{K,\Lambda}\bm w\) both lie in the normal space
\[
    \N_{K,\Lambda,T}^{\perp_T}.
\]
Inference for \(\beta\) is therefore carried by the projected moment
\[
    \Psi_T(\beta,\Lambda)
    =
    \innerT{\bm M_{K,\Lambda}(\bm y-\beta\bm x)}
            {\bm M_{K,\Lambda}\bm w}.
\]

\section{Three Bridge Lemmas}

The following three finite-sample facts are the minimum material needed before
the good-event proof.

\begin{lemma}[Profiling is distance to a nuisance subspace]
\label{lem:min-rss-distance}
For fixed \((\beta,\Lambda)\),
\[
    RSS_K(\beta,\Lambda)
    =
    T\operatorname{dist}_T^2(\bm y-\beta\bm x,\N_{K,\Lambda,T})
    =
    T\|\bm M_{K,\Lambda}(\bm y-\beta\bm x)\|_T^2.
\]
\end{lemma}

\begin{proof}
The profile RSS minimizes over a separate sieve coefficient vector \(c_j\) on
each segment:
\[
    RSS_K(\beta,\Lambda)
    =
    \min_{c_0,\ldots,c_q}
    \sum_{j=0}^q\sum_{t\in I_j(\Lambda)}
    \{Y_t-\beta X_{t-1}-P_K(W_{t-1})'c_j\}^2.
\]
Stacking \(c=(c_0',\ldots,c_q')'\), the fitted nuisance vector is
\(\bm Q_{K,\Lambda}c\). Hence
\[
    RSS_K(\beta,\Lambda)
    =
    \min_c\|\bm y-\beta\bm x-\bm Q_{K,\Lambda}c\|_2^2
    =
    \inf_{g\in\N_{K,\Lambda,T}}\|\bm y_\beta-g\|_2^2.
\]
Since \(\|a\|_2^2=T\|a\|_T^2\), and since the orthogonal projection gives the
nearest point in \(\N_{K,\Lambda,T}\), the claim follows.
\end{proof}

\begin{lemma}[Projected moment and population identification]
\label{lem:min-pop-ident}
Let the true population nuisance space be
\[
    \Hh_0
    =
    \left\{
    h(A)=\sum_{j=0}^{q_0}h_j(W)
    \ind\{\tau_{j,0}<S\le \tau_{j+1,0}\}:h_j\in L^2(P_W)
    \right\},
\]
where \(A=(S,W)\). Let \(M_0^P\) be the population residual-maker onto
\(\Hh_0^\perp\). If \(\mu_0\in\Hh_0\), \(\E(U\mid X,A)=0\), and
\[
    \Gamma=\innerP{M_0^PX}{M_0^Pw}\ne0,
\]
then
\[
    \Psi(\beta,\Lambda_0)
    =
    \innerP{M_0^P(Y-\beta X)}{M_0^Pw}
    =
    -(\beta-\beta_0)\Gamma.
\]
Therefore \(\beta_0\) is the unique population zero.
\end{lemma}

\begin{proof}
At the truth,
\[
    Y-\beta X=\mu_0+U-(\beta-\beta_0)X.
\]
Because \(\mu_0\in\Hh_0\), \(M_0^P\mu_0=0\). Because \(M_0^Pw\) is a function
orthogonal to the nuisance space and \(\E(U\mid X,A)=0\), the population
exogeneity condition gives
\[
    \innerP{M_0^PU}{M_0^Pw}=0.
\]
Thus
\[
    \Psi(\beta,\Lambda_0)
    =
    -(\beta-\beta_0)\innerP{M_0^PX}{M_0^Pw}
    =
    -(\beta-\beta_0)\Gamma.
\]
If \(\Gamma\ne0\), this equals zero only when \(\beta=\beta_0\).
\end{proof}

\begin{lemma}[Scalar EL is convex KKT]
\label{lem:min-el-kkt}
For fixed \((\beta,\Lambda)\), define
\[
    g_t(\beta,\Lambda)
    =
    \{\bm M_{K,\Lambda}(\bm y-\beta\bm x)\}_t
    \{\bm M_{K,\Lambda}\bm w\}_t.
\]
The one-dimensional empirical likelihood problem is
\[
    \max_{\pi\in\Delta_T}\sum_{t=1}^T\log(T\pi_t)
    \quad\text{subject to}\quad
    \sum_{t=1}^T\pi_tg_t(\beta,\Lambda)=0.
\]
If zero lies in the interior of the convex hull of \(\{g_t\}_{t=1}^T\), the KKT
weights have the form
\[
    \pi_t(\lambda)=\frac{1}{T\{1+\lambda g_t\}},
    \qquad
    P_T\left\{\frac{g_t}{1+\lambda g_t}\right\}=0,
\]
where \(P_Ta_t=T^{-1}\sum_ta_t\). If additionally
\[
    P_Tg_t=O(T^{-1/2}),\qquad
    P_Tg_t^2\ge c>0,\qquad
    \max_t|g_t|/\sqrt T=o(1),
\]
then
\[
    -2\log EL(\beta,\Lambda)
    =
    \frac{T(P_Tg_t)^2}{P_Tg_t^2}+o(1).
\]
\end{lemma}

\begin{proof}
The feasible set is the intersection of the probability simplex with the affine
moment hyperplane \(\sum_t\pi_tg_t=0\). The Lagrangian first-order condition for
\(\pi_t\), together with \(\sum_t\pi_t=1\), gives
\[
    \pi_t=\frac{1}{T(1+\lambda g_t)}
\]
for a scalar multiplier \(\lambda\). The moment restriction gives the displayed
dual equation. Under the size conditions, the dual equation implies
\[
    \lambda=\frac{P_Tg_t}{P_Tg_t^2}+o(T^{-1/2})
\]
and \(\max_t|\lambda g_t|=o(1)\). Expanding
\[
    -2\log EL(\beta,\Lambda)
    =
    2\sum_{t=1}^T\log(1+\lambda g_t)
\]
to second order gives the quadratic formula.
\end{proof}

\section{The Good Event}

The proof separates deterministic geometry from probability. Let
\(\delta_T\downarrow0\). The good event is
\[
    \mathcal G_T
    =
    G_T^{\mathrm{iso}}
    \cap G_T^{\mathrm{approx}}
    \cap G_T^{\mathrm{angle}}
    \cap G_T^{\mathrm{local}}
    \cap G_T^{\mathrm{select}}
    \cap G_T^{\mathrm{convex}}.
\]
The components mean:
\begin{enumerate}[leftmargin=2em]
\item \(G_T^{\mathrm{iso}}\): empirical Gram matrices and the empirical inner
products used in the score, derivative, and variance are stable on the searched
subspaces.
\item \(G_T^{\mathrm{approx}}\): the true nuisance is close to the true
break-aware sieve subspace. With \(\bm M_0=\bm M_{K,\Lambda_0}\),
\(\bm r_{\mu,T}=\bm M_0\bm\mu_0\) satisfies
\[
    \|\bm r_{\mu,T}\|_T\le\delta_T,\qquad
    T^{-1/2}|\bm r_{\mu,T}'\bm M_0\bm w|\le\delta_T.
\]
\item \(G_T^{\mathrm{angle}}\): the population transversality constant
\(|\Gamma|\) is bounded away from zero, and the oracle score variance is
nondegenerate. With \(\bm v_0=\bm M_0\bm w\),
\[
    A_T^0=T^{-1/2}\bm u'\bm v_0,\qquad
    B_T^0=T^{-1}\sum_{t=1}^Tu_t^2v_{0,t}^2,
\]
we have \(|A_T^0|\le C\) and \(B_T^0\ge c>0\).
\item \(G_T^{\mathrm{local}}\): the estimated partition can replace the true
partition in the only directions needed by the score:
\[
    \sqrt T|\Psi_T(\beta_0,\widehat\Lambda)-\Psi_T(\beta_0,\Lambda_0)|
    \le\delta_T,
\]
\[
    |V_T(\widehat\Lambda)-V_T(\Lambda_0)|\le\delta_T,
    \qquad
    \sup_{\beta\in\mathcal B}
    |\Psi_T(\beta,\widehat\Lambda)-\Psi_T(\beta,\Lambda_0)|\le\delta_T.
\]
Equivalently, residual-makers are close in the directional topology actually
used by the proof, not necessarily in global operator norm.
\item \(G_T^{\mathrm{select}}\): RSS selects a local break partition and, under
unknown order, the penalized criterion selects \(q_0\).
\item \(G_T^{\mathrm{convex}}\): the scalar EL moment cloud straddles zero and
\(\max_t|g_t(\beta_0,\widehat\Lambda)|/\sqrt T\le\delta_T\), so the EL KKT root
exists and the logarithmic barrier expansion is valid.
\end{enumerate}

\begin{center}
\begin{tabular}{p{0.27\textwidth}p{0.6\textwidth}}
\toprule
Component & Used for \\
\midrule
\(G_T^{\mathrm{iso}}\) & stable projection geometry and denominator control \\
\(G_T^{\mathrm{approx}}\) & remove the nuisance from the oracle score \\
\(G_T^{\mathrm{angle}}\) & identification, nonzero derivative, nonzero variance \\
\(G_T^{\mathrm{local}}\) & replace \(\Lambda_0\) by \(\widehat\Lambda\) \\
\(G_T^{\mathrm{select}}\) & justify break localization and order selection \\
\(G_T^{\mathrm{convex}}\) & justify EL KKT existence and quadratic expansion \\
\bottomrule
\end{tabular}
\end{center}

The full-detail version proves
\[
    \Pp(\mathcal G_T)\to1.
\]
The main body uses only the deterministic implications of \(\mathcal G_T\).

\section{Deterministic Geometry on the Good Event}

\begin{theorem}[Consistency on \(\mathcal G_T\)]
\label{thm:min-consistency}
Suppose Lemma~\ref{lem:min-pop-ident} holds. On \(\mathcal G_T\), if
\[
    |\Psi_T(\widehat\beta,\widehat\Lambda)|\le\delta_T,
\]
then
\[
    |\widehat\beta-\beta_0|\le \frac{2\delta_T}{|\Gamma|}.
\]
Consequently, if \(\Pp(\mathcal G_T)\to1\), then
\(\widehat\beta\to_p\beta_0\).
\end{theorem}

\begin{proof}
On \(\mathcal G_T\), the uniform bridge between sample and population moments
gives
\[
    \sup_{\beta\in\mathcal B}
    |\Psi_T(\beta,\widehat\Lambda)-\Psi(\beta,\Lambda_0)|\le\delta_T.
\]
Therefore
\[
    |\Psi(\widehat\beta,\Lambda_0)|
    \le
    |\Psi(\widehat\beta,\Lambda_0)-\Psi_T(\widehat\beta,\widehat\Lambda)|
    +|\Psi_T(\widehat\beta,\widehat\Lambda)|
    \le 2\delta_T.
\]
By Lemma~\ref{lem:min-pop-ident},
\[
    |\Psi(\widehat\beta,\Lambda_0)|
    =
    |\widehat\beta-\beta_0||\Gamma|,
\]
so the claim follows.
\end{proof}

\begin{theorem}[Wilks reduction on \(\mathcal G_T\)]
\label{thm:min-wilks-reduction}
Let
\[
    V_T(\Lambda)=P_Tg_t(\beta_0,\Lambda)^2.
\]
On \(\mathcal G_T\),
\[
    -2\log EL(\beta_0,\widehat\Lambda)
    =
    \frac{(A_T^0)^2}{B_T^0}+o(1),
\]
where
\[
    A_T^0=T^{-1/2}\bm u'\bm v_0,
    \qquad
    B_T^0=T^{-1}\sum_{t=1}^Tu_t^2v_{0,t}^2,
    \qquad
    \bm v_0=\bm M_0\bm w.
\]
If additionally
\[
    A_T^0\Rightarrow N(0,V),
    \qquad
    B_T^0\to_p V>0,
\]
then
\[
    -2\log EL(\beta_0,\widehat\Lambda)\Rightarrow \chi_1^2.
\]
\end{theorem}

\begin{proof}
By Lemma~\ref{lem:min-el-kkt} and \(G_T^{\mathrm{convex}}\),
\[
    -2\log EL(\beta_0,\widehat\Lambda)
    =
    \frac{T\{\Psi_T(\beta_0,\widehat\Lambda)\}^2}
         {V_T(\widehat\Lambda)}+o(1).
\]
By \(G_T^{\mathrm{local}}\), \(\widehat\Lambda\) can be replaced by
\(\Lambda_0\) in both numerator and denominator. At the true partition,
\[
    \sqrt T\Psi_T(\beta_0,\Lambda_0)
    =
    T^{-1/2}\bm u'\bm v_0
    +
    T^{-1/2}\bm r_{\mu,T}'\bm v_0
    =
    A_T^0+o(1),
\]
where the last equality uses \(G_T^{\mathrm{approx}}\). Also
\[
    V_T(\Lambda_0)=B_T^0+o(1)
\]
by \(G_T^{\mathrm{iso}}\) and the variance part of \(G_T^{\mathrm{angle}}\).
Substitution gives the good-event reduction. The chi-square limit follows by
Slutsky.
\end{proof}

\begin{proposition}[Break search is subspace selection]
\label{prop:min-break-search}
By Lemma~\ref{lem:min-rss-distance}, the break estimator chooses the nuisance
subspace closest to \(\bm y-\beta\bm x\):
\[
    \widehat\Lambda_q(\beta)
    \in
    \arg\min_{\Lambda\in\mathcal L_q(\epsilon)}
    \operatorname{dist}_T^2(\bm y-\beta\bm x,\N_{K,\Lambda,T}).
\]
On \(G_T^{\mathrm{select}}\), this nearest-subspace selector is local around
\(\Lambda_0\), and the penalized selector chooses \(q_0\).
\end{proposition}

\section{Main Theorem}

\begin{theorem}[Profile-sieve EL main theorem]
\label{thm:min-main}
Assume:
\begin{enumerate}[leftmargin=2em]
\item the model and population identification conditions in
Lemma~\ref{lem:min-pop-ident};
\item the stochastic regularity conditions imply \(\Pp(\mathcal G_T)\to1\);
\item the oracle projected score satisfies
\[
    A_T^0=T^{-1/2}\bm u'\bm M_0\bm w\Rightarrow N(0,V),
    \qquad
    B_T^0=T^{-1}\sum_{t=1}^Tu_t^2(\bm M_0\bm w)_t^2\to_p V>0;
\]
\item the estimator satisfies the local projected moment condition, with
probability tending to one,
\[
    |\Psi_T(\widehat\beta,\widehat\Lambda)|\le \delta_T.
\]
\end{enumerate}
Then
\[
    \widehat\beta\to_p\beta_0
\]
and
\[
    -2\log EL(\beta_0,\widehat\Lambda)\Rightarrow \chi_1^2.
\]
\end{theorem}

\begin{proof}
Consistency follows from Theorem~\ref{thm:min-consistency} because
\(\Pp(\mathcal G_T)\to1\). The Wilks result follows from
Theorem~\ref{thm:min-wilks-reduction}, the oracle score CLT, and the oracle
variance LLN.
\end{proof}

\section{Efficiency Add-On}

Efficiency is not needed for Wilks validity. It requires an additional
efficiency model. Let \(A=(S,W)\), and suppose the conditional Gaussian location
calculation is valid:
\[
    U\mid X,A\sim N(0,\sigma^2(A)),
    \qquad 0<c\le\sigma^2(A)\le C<\infty.
\]
The raw score is \(UX/\sigma^2(A)\), and the nuisance tangent directions are
\[
    Uh(A)/\sigma^2(A),\qquad h\in\Hh_0.
\]
Projecting the raw score away from the nuisance tangent space gives
\[
    S_{\mathrm{eff}}
    =
    \frac{U\{X-E(X\mid A)\}}{\sigma^2(A)}.
\]
The efficient information is
\[
    I_{\mathrm{eff}}
    =
    \E\left[
    \frac{\{X-E(X\mid A)\}^2}{\sigma^2(A)}
    \right].
\]
To make the projected sample score equal this efficient score, choose
\[
    w^*(X,A)=\frac{X}{\sigma^2(A)}.
\]
Because \(\sigma^2(A)\) is \(A\)-measurable,
\[
    M_0^Pw^*
    =
    \frac{X-E(X\mid A)}{\sigma^2(A)}.
\]
Thus, at \(\beta_0\),
\[
    \{M_0^P(Y-\beta_0X)\}\{M_0^Pw^*\}
    =
    \frac{U\{X-E(X\mid A)\}}{\sigma^2(A)}
    =
    S_{\mathrm{eff}}.
\]
If the sieve and estimated-break replacements preserve this score at
\(T^{-1/2}\) scale, then any EL or moment root using \(w^*\) has
\[
    \sqrt T(\widehat\beta-\beta_0)
    =
    I_{\mathrm{eff}}^{-1}T^{-1/2}\sum_{t=1}^T
    \frac{U_t\{X_t-E(X_t\mid A_t)\}}{\sigma^2(A_t)}
    +o_p(1),
\]
and hence attains the semiparametric efficiency bound
\[
    I_{\mathrm{eff}}^{-1}
    =
    \left\{
    \E\left[
    \frac{\{X-E(X\mid A)\}^2}{\sigma^2(A)}
    \right]
    \right\}^{-1}.
\]
In the homoskedastic case, \(w^*=X/\sigma^2\), so the same root is obtained by
using \(w=X\).

\section{What This Minimal Body Leaves to the Full Version}

The minimal body does not reprove the probability bounds. The full-detail body
contains the supporting proofs for:
\begin{enumerate}[leftmargin=2em]
\item empirical Gram stability and empirical-to-population projection bridges;
\item sieve approximation of the true nuisance and projected score directions;
\item RSS localization and unknown-order penalty comparisons;
\item estimated-break directional replacement;
\item oracle score CLT and variance LLN;
\item max-score and convex-hull control for EL;
\item heteroskedastic efficient-weight replacement, if the efficiency theorem is
used.
\end{enumerate}
The role of this minimal body is to show exactly how those stochastic facts are
used once they are available.

\end{document}
